\documentclass[11pt,letterpaper]{article}
\usepackage[T1]{fontenc}
\usepackage{lmodern}
\usepackage[hmargin=1.05in,vmargin=1in]{geometry}
\usepackage{amsmath,amssymb,amsthm}
\usepackage{microtype}
\usepackage{xcolor}
\definecolor{linkblue}{RGB}{30,69,101}
\usepackage[colorlinks=true,linkcolor=linkblue,citecolor=linkblue,urlcolor=linkblue]{hyperref}
\hypersetup{
  pdftitle={A superpolynomial lower bound for fixed-flow continuation in Bayesian simple security games},
  pdfauthor={Robert Thomas Collins}
}
\theoremstyle{definition}
\newtheorem{theorem}{Theorem}
\newtheorem{lemma}[theorem]{Lemma}
\newcommand{\X}{\mathcal{X}}
\newcommand{\ind}{\mathbf{1}}
\newcommand{\recordlink}[2]{\href{https://github.com/robbiemu/bayesian-simple-security-games/blob/fd0e76a3090ffaf4b3f85cd166fb57b69285d3be/#1}{#2}}
\setlength{\emergencystretch}{1.5em}
\raggedbottom

\begin{document}
\begin{center}
  {\LARGE\bfseries A superpolynomial lower bound for fixed-flow\\[.10em]
    continuation in Bayesian simple security games\par}
  \medskip{\large Robert Thomas Collins\par}
  \medskip{\small Working draft \enspace\textperiodcentered\enspace 21 November 2026\par}
\end{center}

\begin{abstract}
One way to find an equilibrium in a security game is to start with no defender
coverage and increase it while the attacker keeps playing a best response.
How often must the attack distribution change along the way? We construct
one-resource games in which every such continuation that holds attack flow
fixed during each budget increase needs \(L^{\Omega(\log L)}\) segments,
where \(L\) is the binary input length. Distinct shortest paths in a parametric
graph prescribe which target each type attacks on successive budget intervals.
Every equilibrium obeys the prescription, even as total attack probability
stays equal across targets. On the same games an assignment and a linear
program compute an equilibrium at the final budget in polynomial time.
\end{abstract}

\section{Following equilibria as coverage increases}

A defender has one resource and must choose which target to cover. Only
the attacker observes its type before choosing a target. Both players may
randomize. Increasing the probability of covering a target makes it less
attractive to an attacker. Eventually a type may have to move its attack
elsewhere for the two players to remain in equilibrium.

Li, Conitzer and Korzhyk use this observation in their algorithm for
catcher-evader games, which include Bayesian simple security games. They
start with zero defender resources and increase coverage in steps, reallocating
attacks between steps. They ask whether their algorithm always runs in
polynomial time~\cite[\S\S3.1, 6.2 and 7]{LCK}. Here a single completed run
can require superpolynomially many steps, whichever eligible moves the
algorithm chooses.

Write \(x_v\) for the probability of covering target \(v\) and \(b=\sum_v x_v\)
for the current budget. At an intermediate budget the defender can cover
one target with probability \(b\) and leave the resource idle with probability
\(1-b\). At \(b=1\) this is the one-resource game. The \emph{attack flow}
\(q_{iv}\) is the joint probability of type \(i\) and an attack on \(v\).
A \emph{completed fixed-flow continuation} follows equilibria in increasing
budget order from \(0\) to \(1\). Its positive-length budget segments,
together with their endpoints, cover \([0,1]\). On each segment the entire
matrix \(q\) stays fixed and the players are in equilibrium at every
intervening budget. Attacks may be rerouted at a shared endpoint. Segments
may have any length.

\begin{theorem}\label{thm:main}
There is an explicit family of rational Bayesian simple security games with
\(n\) targets, \(n\) equally likely attacker types and one defender resource
in which every completed fixed-flow continuation needs \(n^{\Omega(\log n)}\)
positive-length segments. In the binary length \(L\) of the explicitly
listed game the lower bound is \(L^{\Omega(\log L)}\). Covering the attacked
target raises the defender's payoff by one and lowers the attacker's payoff
by an amount within \(1/(64n)\) of one.

Throughout \([1/4,1]\), every equilibrium gives every target interior
coverage and total attack probability \(1/n\). Within this range there are
\(K=n^{\Omega(\log n)}\) disjoint intervals of positive length. On each
interval a single permutation prescribes the target attacked by each type
in every equilibrium. The \(K\) permutations are distinct.
\end{theorem}

\section{From shortest paths to assignments}\label{sec:paths}

We first construct a table of affine weights for which different permutations
win in turn. Fix an integer \(h\ge4\), put \(m=\lfloor\log_2 h\rfloor\), and
let \(K=h^m\). Take \(B=1\) and \(D=0\) in the parametric shortest-path
construction of Gajjar and Radhakrishnan~\cite[Lemma~14 and proof of
Theorem~1]{GR}. Their layered directed acyclic graph has one source \(o\)
and rational affine arc costs in a parameter \(\lambda\). For \(K\) disjoint
intervals \(I_1,\ldots,I_K\) of positive length, in increasing order, they
give paths \(P_j\) from \(o\) to the last layer. Throughout \(I_j\), every
other path from \(o\) to that layer costs at least one more than \(P_j\)
(property~(iii)). The \(P_j\) are distinct (property~(vi)).

Join every last-layer vertex to a new sink \(t\) by a zero-cost arc. If
arcs are parallel, subdivide each through its own vertex, assigning the
original cost to the first half and zero to the second. Paths to the last
layer now correspond bijectively to \(o\)-to-\(t\) paths of the same cost.
Pad with isolated vertices if needed to bring the vertex count \(n\) up to \(h\).

Use the vertices as both the rows and the columns of an assignment. Allow
each graph arc at its negated cost, a zero-weight loop at each vertex
except \(t\), and a zero-weight return arc \(t\to o\). An allowed
permutation must use the return arc, the only choice in row \(t\).
Deleting it from its cycle leaves an \(o\)-to-\(t\) path. Every other
cycle is a loop by acyclicity. Conversely, closing a path and fixing its
off-path vertices gives exactly one permutation. Its weight is the
negative of the path cost, so the permutation \(\pi_j\) of \(P_j\) beats
every other allowed permutation by at least one on \(I_j\). This is the
path-to-matching correspondence of~\cite[Appendix~B]{GR} in assignment form.

For \(1/4\le b\le1\), mean coverage \(s=b/n\) ranges from \(1/(4n)\) to
\(1/n\). Let \(\lambda_-\) and \(\lambda_+\) be the midpoints of the first
and last intervals. Set
\begin{equation}
  \lambda(s)=\lambda_-+\frac{4ns-1}{3}(\lambda_+-\lambda_-),\qquad
  J_j=\{b\in[1/4,1]\mid\lambda(b/n)\in I_j\}.
  \label{eq:rescale}
\end{equation}
This increasing map takes the two endpoint budgets to the two midpoints.
Even the first and last intervals retain half their length in the
\(\lambda\) coordinate. Thus all \(J_j\) have positive length, and
\(1\in J_K\).

Write each allowed weight after substitution as \(w_{iv}(s)=B_{iv}-C_{iv}s\).
We need a payoff for \emph{every} type--target pair. Put
\[
  H=1+\Bigl\lceil\max_{iv\ \mathrm{allowed}}(|B_{iv}|+|C_{iv}|)\Bigr\rceil,
  \qquad M=(2n+1)H,
\]
and set \(B_{iv}=-M,\ C_{iv}=0\) on missing pairs. Allowed weights lie
in \([-H,H]\) for \(0\le s\le1\). A permutation using a missing pair
therefore weighs at most \(-M+(n-1)H=-(n+2)H\), while every allowed
permutation weighs at least \(-nH\). The full table satisfies
\begin{equation}
  W_\pi(s)=\sum_i w_{i,\pi(i)}(s),\qquad
  W_{\pi_j}(s)-W_\rho(s)\ge1
  \quad (ns\in J_j,\ \rho\ne\pi_j).
  \label{eq:gap}
\end{equation}
Also \(|B_{iv}|+|C_{iv}|\le M\) on every pair.

\section{What the game forces at a fixed budget}\label{sec:game}

Li, Conitzer and Korzhyk encode a fixed matching problem in a catcher-evader
game~\cite[\S6.1]{LCK}. We use a single game to realize the assignments
in~\eqref{eq:gap} as its budget changes. Give each type prior \(1/n\).
The defender earns \(0\) if the attacked target is uncovered and \(1\)
if it is covered. Type \(i\) earns \(a_{iv}\)
and \(a_{iv}-d_{iv}\), respectively, for an attack on \(v\), where
\[
  \varepsilon=\frac1{64nM^2},\qquad
  a_{iv}=\varepsilon B_{iv},\qquad d_{iv}=1+\varepsilon C_{iv}.
\]
Covering a target reduces the attacker's payoff by \(d_{iv}>0\), with
\(|d_{iv}-1|\le\varepsilon M\le1/(64n)\). At coverage \(x\) the utility
of attacking \(v\) is
\[
  u_{iv}(x)=a_{iv}-d_{iv}x_v=-x_v+\varepsilon w_{iv}(x_v).
\]
The first term draws attacks toward the least-covered targets. The second
carries the assignment weights and has magnitude at most \(\varepsilon M\).

At equilibrium the nonnegative flow in row \(i\) sums to \(1/n\) and uses
only targets maximizing \(u_{iv}(x)\). Against the target totals
\(z_v=\sum_iq_{iv}\), the defender maximizes \(\sum_v x_vz_v\) over
\(\X_b=\{x\in[0,1]^n\mid\sum_v x_v=b\}\). An equilibrium exists at each
\(0\le b\le1\). Indeed, \(\X_b\) is the convex hull of the \(n\) allocations
\(b e_v\) that put the whole budget on one target. These defender actions
and the attacker's maps from types to targets form a finite game with a mixed
Nash equilibrium~\cite{Nash}. Payoffs depend only on \(x,q\), and a profitable
deviation by a single type would improve the attacker's payoff since its prior
is positive.

The defender can put the whole budget on a target of maximum mass
\(z_*=\max_v z_v\), earning \(bz_*\). The shortfall at any feasible \(x\) is
\begin{equation}
  bz_*-\sum_v x_vz_v=\sum_v x_v(z_*-z_v).
  \label{eq:defender}
\end{equation}
Each term is nonnegative. Hence a coverage is optimal exactly when every
positively covered target has mass \(z_*\).

At \(b>0\) a most-covered target \(v\) therefore has mass \(z_*\ge1/n\).
Some type \(i\) attacks it, so \(u_{iv}(x)\ge u_{iw}(x)\) for a
least-covered target \(w\). Since \(|u_{it}(x)+x_t|\le\varepsilon M\)
for each target \(t\), we get \(x_v-x_w\le2\varepsilon M\). The mean
\(s=b/n\) lies between the extremes, giving
\begin{equation}
  |x_t-s|\le2\varepsilon M=\frac1{32nM}\qquad\text{for all targets }t.
  \label{eq:spread}
\end{equation}
For \(1/4\le b\le1\), every coverage is at least
\(1/(4n)-1/(32n)>0\). Each is also below one since \(n\ge2\) and the
sum is at most one. By~\eqref{eq:defender} all targets have maximum mass.
Their masses sum to one, so \(z_v=1/n\) for every \(v\).

\begin{lemma}\label{lem:force}
Fix \(1/4\le b\le1\) and let \(s=b/n\). If \(W_\pi(s)\ge W_\rho(s)+1\)
for every permutation \(\rho\ne\pi\), then every equilibrium at \(b\) has
flow \(q=q^\pi\), where \(q^\pi_{iv}=\ind_{\{v=\pi(i)\}}/n\).
\end{lemma}

\begin{proof}
Fix an equilibrium \((x,q)\). The priors give \(q\) row sums \(1/n\).
All coverages are positive, so~\eqref{eq:defender} gives the same column
sums. Thus \(nq\) is doubly stochastic. By the Birkhoff--von Neumann theorem it is a
convex combination of permutation matrices~\cite{Birkhoff}. If \(\rho\) has
coefficient \(\alpha_\rho>0\), then \(nq_{i,\rho(i)}\ge\alpha_\rho\) for every \(i\).
Thus \(\rho(i)\) is a best response for type \(i\). Consequently
\(\rho\) maximizes \(\sum_i u_{i,\rho(i)}(x)\) among all permutations,
even when a type has several best responses.

Every permutation uses each target once, so its common coverage terms sum
to \(-\sum_v x_v=-b\). The remaining utility is
\(\varepsilon\sum_i w_{i,\rho(i)}(x_{\rho(i)})\).
By~\eqref{eq:spread}, the assignment sum at the actual coverages satisfies
\[
 \left|\sum_i w_{i,\rho(i)}(x_{\rho(i)})-W_\rho(s)\right|
 \le \sum_i |C_{i,\rho(i)}|\,|x_{\rho(i)}-s|
 \le 2n\varepsilon M^2=\tfrac1{32}.
\]
The two errors in a comparison cost at most \(1/16\) of the unit gap.
Thus \(\pi\) is the unique maximizer of the summed
utilities. Every positive coefficient in the decomposition belongs to
\(\pi\), which gives \(q=q^\pi\).
\end{proof}

\section{Counting segments in one run}\label{sec:count}

Fix one completed continuation. By~\eqref{eq:gap} and Lemma~\ref{lem:force},
every equilibrium on \(J_j\) has flow \(q^{\pi_j}\). If the run has finitely
many positive-length segments, choose a budget in the interior of each
\(J_j\) away from all segment endpoints. The segments and their endpoints
cover \([0,1]\), so each chosen budget lies inside a segment carrying
\(q^{\pi_j}\). The \(K\) distinct matrices require \(K\) different segments
of this run. Infinitely many positive-length segments already meet the
bound. In either case the prescribed flows occur in order and require
at least \(K-1\) changes.

To express the count in bits, consider the graph before subdivision. With
its sink it has at most \(6h^8+1\) vertices, and its arc coefficients use
\(O((\log h)^3)\) bits per numerator and denominator~\cite[Lemma~14(i)--(ii)
and proof of Theorem~1]{GR}. Its \(O(3^m)\) recursive pieces each use at most
\((1+mh)(h+1)\) straight segments in a linking graph~\cite[\S\S2.1 and 3.3]{GR}.
Each pair crosses at most once, giving \(h^{O(1)}\) arcs after splitting at
crossings. Subdivision and padding therefore give \(h\le n\le h^{O(1)}\).

All interval endpoints are integers in \([0,(h^2)^{m+1}]\). The smallest
and largest digit at each level of the interval recurrence give the first
and last intervals~\cite[\S3.1]{GR}. Their midpoints have denominator at
most two and \(O((\log h)^2)\) bits. Substitution~\eqref{eq:rescale} takes
only a constant number of rational operations per arc, giving
\(O((\log n)^3)\) bits per allowed coefficient \(B_{iv},C_{iv}\). The
maximum and ceiling in \(H\) and the formulas for \(M,\varepsilon\) and
the payoffs preserve this bound per entry.

Listing the \(n^2\) type--target pairs gives \(L=\Omega(n^2)\) and
\(L=O(n^2(\log n)^3)\). The intervals serve only as witnesses. Since
\(\log L=\Theta(\log n)=\Theta(\log h)\), Theorem~\ref{thm:main} follows from
\[
  \log K=\lfloor\log_2 h\rfloor\log h
        =\Theta\bigl((\log n)^2\bigr)=\Theta\bigl((\log L)^2\bigr).
\]

\section{The published continuation and the final budget}

In the reduction of Li, Conitzer and Korzhyk, evader \(i\)'s allocation is
row \(i\) of \(q\). With one attack per type, its resource and its cap at
each target are both \(1/n\), and its payoff change upon coverage is
\(-d_{iv}\)~\cite[\S3.1]{LCK}. On \([1/4,1]\) every target is below its
defender cap and has maximum mass, so all targets belong to their
\emph{open defender boundary}.

Algorithm~1 completes the traversal from zero to the full defender
budget~\cite[Theorem~5]{LCK}. Algorithms~2 and~4 reroute attacks at
fixed coverage. A successful call to Algorithm~3 moves \(x\) along a
straight line with \(q\) fixed and preserves equilibrium~\cite[Lemma~4]{LCK}.
The attack comparisons on its support are affine in \(x\). The defender's
shortfall~\eqref{eq:defender} is linear in \(x\) since \(b=\sum_v x_v\).
The comparisons hold and the shortfall vanishes at both endpoints, hence
throughout the step. Each successful call is a fixed-flow segment, so
Algorithm~1 requires at least \(K\) such calls for every legal seed choice
and min-cost tie rule. Taking the greatest feasible increase at each fixed
flow also requires \(K\) segments or more.

\paragraph{Reallocation cost.}
Algorithm~2 uses the cost \(\ell(q)=\sum_{iv}q_{iv}\log d_{iv}\), with
natural logarithms. The prescribed flows have strictly decreasing cost in
budget order. Indeed, let \(ns_1<ns_2\) lie in different witness intervals,
with winners \(\pi,\rho\), respectively. Write \(C_\pi=\sum_i C_{i,\pi(i)}\).
The difference \(W_\pi-W_\rho\) is at least \(1\) at \(s_1\) and at
most \(-1\) at \(s_2\). Thus \((C_\pi-C_\rho)(s_2-s_1)\ge2\), giving
\(C_\pi-C_\rho\ge8n/3\) since \(s_2-s_1\le3/(4n)\).
For \(|t|\le1/2\), integration of \(|t/(1+t)|\le2|t|\) gives
\(|\log(1+t)-t|\le t^2\). Applying this with \(t=\varepsilon C_{iv}\) gives
\[
  n\bigl(\ell(q^\pi)-\ell(q^\rho)\bigr)
  \ge\varepsilon(C_\pi-C_\rho)-2n\varepsilon^2M^2
  \ge\varepsilon(8n/3-1/32)>0 .
\]

\paragraph{Computing the endpoint directly.}
Conitzer asked whether one Bayes--Nash equilibrium of a given game can be
found in polynomial time~\cite{Conitzer}. On this family two rational linear
programs suffice at \(b=1\). First maximize the assignment objective with weights
\(a_{iv}-d_{iv}/n=-1/n+\varepsilon w_{iv}(1/n)\) over the doubly stochastic
matrices. Since \(1\in J_K\), the unique optimum is the permutation
matrix of \(\pi=\pi_K\). Next find any \(x\in\X_1\) satisfying
\[
  a_{i,\pi(i)}-d_{i,\pi(i)}x_{\pi(i)}
    \ge a_{iv}-d_{iv}x_v\qquad\text{for all }i,v .
\]
Existence of an equilibrium with flow \(q^\pi\) guarantees feasibility.
Every solution makes the assigned attacks optimal, and their uniform
column totals make every \(x\in\X_1\) optimal for the defender. Both
programs have polynomial encoding length and can be solved in polynomial
bit time~\cite{Khachiyan}. The defender covers target \(v\) with probability \(x_v\).

\section*{Contributions and use of AI}

Robert Thomas Collins asked about the longest single run in each class of
games, motivated by the Busy Beaver problem. He directed the investigation
and reviewed the scope and presentation of the result. In OpenAI Codex,
GPT-6.2 Sol reviewed sources and tested small examples in exact arithmetic.
GPT-6.2 Astra helped derive and check the reduction from parametric paths
to forced attack flows and drafted the proof. Gajjar and Radhakrishnan
proved the shortest-path result, including the interval-wide gaps~\cite{GR}.

A separate model-assisted audit checked the reduction, the segment count,
the input length and the endpoint argument. In the four-target example,
enumerating all \(24\) assignments gives three different winners at the
three reported budgets, each with gap two. In the three-target audit,
testing the row inequalities for all six permutations gives one feasible
assignment at each of two budgets. The project record supplies both
examples~\cite{record}.

The Lean checks accompanying this edition cover affine path sums,
permutation sums and the decomposition of a given doubly stochastic
matrix. They also check that an open interval contains a point outside any
given finite set and that an injection between finite sets bounds their
sizes. Run \texttt{lake build LibraryChecks} from the repository's
\texttt{lean/} directory to reproduce them~\cite{record}. The forcing
lemma, the segment count and the binary input-length bound are verified
by the written argument in Sections~\ref{sec:paths}--\ref{sec:count}.

\begin{thebibliography}{9}
\small
\raggedright
\setlength{\itemsep}{.25em}
\bibitem{LCK}
Y. Li, V. Conitzer, and D. Korzhyk.
\emph{Catcher-Evader Games}. IJCAI, 2016.
\href{https://arxiv.org/pdf/1602.01896v2}{arXiv:1602.01896v2}.

\bibitem{GR}
K. Gajjar and J. Radhakrishnan.
\emph{Parametric Shortest Paths in Planar Graphs}. FOCS, 2019.
\href{https://arxiv.org/pdf/1811.05115v3}{arXiv:1811.05115v3}.

\bibitem{Nash}
J. F. Nash, Jr. \href{https://doi.org/10.1073/pnas.36.1.48}{\emph{Equilibrium
points in \(n\)-person games}}. Proc. Natl. Acad. Sci. USA 36, 48--49, 1950.

\bibitem{Birkhoff}
G. Birkhoff. \emph{Tres observaciones sobre el \'algebra lineal}.
Univ. Nac. Tucum\'an, Ser. A 5, 147--151, 1946.

\bibitem{Conitzer}
V. Conitzer. \emph{Open problem: Computing Bayes-Nash equilibrium in
simple security games}.
\href{https://arxiv.org/pdf/2610.08796v2}{arXiv:2610.08796v2}, 2026.

\bibitem{Khachiyan}
L. G. Khachiyan. \href{https://www.mathnet.ru/eng/zvmmf5239}{\emph{Polynomial
algorithms in linear programming}}. USSR Comput. Math. Math. Phys. 20(1), 53--72, 1980.

\bibitem{record}
R. T. Collins and contributors. \emph{Bayesian simple security games},
project records at revision \texttt{fd0e76a3}.
\recordlink{notes/B3f-forced-flow-switches.md}{Construction and four-target example},
\recordlink{notes/B3g-lower-bound-audit.md}{audit and three-target example},
\recordlink{lean/README.md}{Lean checks}, and
\recordlink{records/21a3ca77-54c5-4e22-8493-9926749662d4.md}{contribution record}.
\end{thebibliography}
\end{document}
